\documentclass[aspectratio=169,8pt]{beamer}
\usetheme{Madrid}
\usepackage[utf8]{inputenc}
\usepackage[T1]{fontenc}
\usepackage{amsmath,amssymb}
\usepackage{booktabs}
\usepackage{enumitem}
\setlist[enumerate]{label=\Alph*.,leftmargin=*,itemsep=1pt,topsep=2pt}
\setlist[itemize]{leftmargin=*,itemsep=1pt,topsep=2pt}
\setbeamerfont{block body}{size=\tiny}
\setbeamerfont{block title}{size=\scriptsize}
\setbeamerfont{frametitle}{size=\footnotesize}
\setbeamerfont{normal text}{size=\tiny}
\setbeamertemplate{navigation symbols}{}
\setbeamertemplate{itemize item}{\tiny$\bullet$}
\setbeamertemplate{itemize subitem}{\tiny$\circ$}

\title[GARP RAI]{GARP Risk \& AI --- Numerical Problems (Conceptual Topics)}
\subtitle{Unofficial Exam Prep --- With Detailed Definitions}
\author{Unofficial}
\date{\today}

\begin{document}
	
	\begin{frame}
		\titlepage
	\end{frame}
	
	% =================================================================
	\section{Group Fairness Measures --- Numerical}
	% =================================================================
	
	\begin{frame}{C001 --- Demographic Parity}
		\begin{block}{Question}
			Group A approval rate = 0.70. Group B approval rate = 0.50. What is the demographic parity difference?
			\begin{enumerate}
				\item 0.10
				\item 0.20
				\item 0.30
				\item 0.40
			\end{enumerate}
		\end{block}
		\begin{exampleblock}{Answer: B}\end{exampleblock}
		\begin{block}{Explanation}
			\textbf{Definition:} Demographic parity requires equal positive prediction rates across groups. The difference measures disparity.\\
			\textbf{Formula:} Difference = $|$Rate A $-$ Rate B$|$\\
			\textbf{Calculation:} $|$0.70 $-$ 0.50$|$ = 0.20\\
			\textbf{Interpretation:} There is a 20 percentage point disparity between groups.\\
			\textbf{Learning Objective:} Describe various measures of group fairness.\\
			\textbf{Reference:} Module 3, Chapter 2
		\end{block}
	\end{frame}
	
	\begin{frame}{C002 --- Predictive Rate Parity}
		\begin{block}{Question}
			Group A precision = 0.80. Group B precision = 0.60. What is the predictive rate parity difference?
			\begin{enumerate}
				\item 0.10
				\item 0.20
				\item 0.30
				\item 0.40
			\end{enumerate}
		\end{block}
		\begin{exampleblock}{Answer: B}\end{exampleblock}
		\begin{block}{Explanation}
			\textbf{Definition:} Predictive rate parity requires equal precision across groups.\\
			\textbf{Formula:} Difference = $|$Precision A $-$ Precision B$|$\\
			\textbf{Calculation:} $|$0.80 $-$ 0.60$|$ = 0.20\\
			\textbf{Interpretation:} There is a 20 percentage point difference in precision.\\
			\textbf{Learning Objective:} Describe various measures of group fairness.\\
			\textbf{Reference:} Module 3, Chapter 2
		\end{block}
	\end{frame}
	
	\begin{frame}{C003 --- Equal Opportunity}
		\begin{block}{Question}
			Group A recall = 0.90. Group B recall = 0.70. What is the equal opportunity difference?
			\begin{enumerate}
				\item 0.10
				\item 0.20
				\item 0.30
				\item 0.40
			\end{enumerate}
		\end{block}
		\begin{exampleblock}{Answer: B}\end{exampleblock}
		\begin{block}{Explanation}
			\textbf{Definition:} Equal opportunity requires equal true positive rates (recall) across groups.\\
			\textbf{Formula:} Difference = $|$Recall A $-$ Recall B$|$\\
			\textbf{Calculation:} $|$0.90 $-$ 0.70$|$ = 0.20\\
			\textbf{Interpretation:} There is a 20 percentage point difference in recall.\\
			\textbf{Learning Objective:} Describe various measures of group fairness.\\
			\textbf{Reference:} Module 3, Chapter 2
		\end{block}
	\end{frame}
	
	\begin{frame}{C004 --- Disparate Impact Ratio}
		\begin{block}{Question}
			Group A approval rate = 0.60. Group B approval rate = 0.30. What is the disparate impact ratio?
			\begin{enumerate}
				\item 0.30
				\item 0.50
				\item 0.60
				\item 2.00
			\end{enumerate}
		\end{block}
		\begin{exampleblock}{Answer: B}\end{exampleblock}
		\begin{block}{Explanation}
			\textbf{Definition:} Disparate impact ratio compares selection rates between groups. A ratio below 0.80 may indicate adverse impact.\\
			\textbf{Formula:} Ratio = Protected rate / Reference rate\\
			\textbf{Calculation:} 0.30 / 0.60 = 0.50\\
			\textbf{Interpretation:} The ratio is 0.50, below the 0.80 threshold, indicating potential adverse impact.\\
			\textbf{Learning Objective:} Describe sources of algorithmic bias and unfairness.\\
			\textbf{Reference:} Module 3, Chapter 2
		\end{block}
	\end{frame}
	
	\begin{frame}{C005 --- Sample Representation}
		\begin{block}{Question}
			Dataset has 2000 samples. 300 are from a protected group. What percentage is protected?
			\begin{enumerate}
				\item 10 percent
				\item 15 percent
				\item 20 percent
				\item 25 percent
			\end{enumerate}
		\end{block}
		\begin{exampleblock}{Answer: B}\end{exampleblock}
		\begin{block}{Explanation}
			\textbf{Definition:} Representation measures the proportion of a group in the dataset. Underrepresentation can lead to bias.\\
			\textbf{Formula:} Percentage = Protected / Total $\times$ 100\\
			\textbf{Calculation:} 300/2000 = 0.15 = 15 percent\\
			\textbf{Interpretation:} The protected group represents 15 percent of the dataset.\\
			\textbf{Learning Objective:} Describe sources of algorithmic bias and unfairness.\\
			\textbf{Reference:} Module 3, Chapter 2
		\end{block}
	\end{frame}
	
	\begin{frame}{C006 --- Fairness Impossibility}
		\begin{block}{Question}
			Group A base rate = 0.30. Group B base rate = 0.10. Can all three fairness measures (demographic parity, predictive rate parity, equal opportunity) be satisfied simultaneously?
			\begin{enumerate}
				\item Yes
				\item No
				\item Only if sample size is large
				\item Only if base rates are equal
			\end{enumerate}
		\end{block}
		\begin{exampleblock}{Answer: B}\end{exampleblock}
		\begin{block}{Explanation}
			\textbf{Definition:} When base rates differ between groups, it is mathematically impossible to satisfy all three fairness measures simultaneously (impossibility theorem).\\
			\textbf{Formula:} If base rates are unequal, no model can satisfy demographic parity, predictive rate parity, and equal opportunity at the same time.\\
			\textbf{Calculation:} Base rates 0.30 and 0.10 are unequal, so impossibility applies.\\
			\textbf{Interpretation:} Trade-offs between fairness measures are unavoidable.\\
			\textbf{Learning Objective:} Discuss trade-offs associated with different concepts and measures of fairness.\\
			\textbf{Reference:} Module 3, Chapter 2
		\end{block}
	\end{frame}
	
	% =================================================================
	\section{Naive Bayes --- Numerical}
	% =================================================================
	
	\begin{frame}{C007 --- Naive Bayes Unnormalized Posterior}
		\begin{block}{Question}
			Prior P(Good) = 0.6. Likelihoods: P(w1$|$Good) = 0.5, P(w2$|$Good) = 0.4. What is the unnormalized posterior?
			\begin{enumerate}
				\item 0.06
				\item 0.12
				\item 0.24
				\item 0.40
			\end{enumerate}
		\end{block}
		\begin{exampleblock}{Answer: B}\end{exampleblock}
		\begin{block}{Explanation}
			\textbf{Definition:} Naive Bayes calculates the probability of a class given observed features, assuming features are conditionally independent.\\
			\textbf{Formula:} Unnormalized Posterior = P(Class) $\times$ $\prod$ P(word$_{i}$$|$Class)\\
			\textbf{Calculation:} 0.5 $\times$ 0.4 $\times$ 0.6 = 0.12\\
			\textbf{Interpretation:} The unnormalized score for ``Good'' is 0.12.\\
			\textbf{Learning Objective:} Explain how the Naive Bayes classifier is used to categorize documents.\\
			\textbf{Reference:} Module 2, Chapter 9
		\end{block}
	\end{frame}
	
	\begin{frame}{C008 --- Naive Bayes Normalization}
		\begin{block}{Question}
			Unnormalized posteriors: Good = 0.12, Bad = 0.08. What is the normalized P(Good)?
			\begin{enumerate}
				\item 0.40
				\item 0.50
				\item 0.60
				\item 0.67
			\end{enumerate}
		\end{block}
		\begin{exampleblock}{Answer: C}\end{exampleblock}
		\begin{block}{Explanation}
			\textbf{Definition:} Normalization converts unnormalized posterior scores into proper probabilities that sum to 1.\\
			\textbf{Formula:} P(Good) = Unnorm(Good) / (Unnorm(Good) + Unnorm(Bad))\\
			\textbf{Calculation:} 0.12 / (0.12 + 0.08) = 0.12 / 0.20 = 0.60\\
			\textbf{Interpretation:} Given the observed words, there is a 60 percent probability the document belongs to the ``Good'' class.\\
			\textbf{Learning Objective:} Explain how the Naive Bayes classifier is used to categorize documents.\\
			\textbf{Reference:} Module 2, Chapter 9
		\end{block}
	\end{frame}
	
	\begin{frame}{C009 --- Naive Bayes Smoothing}
		\begin{block}{Question}
			Word count = 0 in class Good. Total words in Good = 60. Vocabulary = 200. Using Laplace smoothing ($\lambda$ = 1), what is the smoothed probability?
			\begin{enumerate}
				\item 0.00
				\item 0.004
				\item 0.01
				\item 0.02
			\end{enumerate}
		\end{block}
		\begin{exampleblock}{Answer: B}\end{exampleblock}
		\begin{block}{Explanation}
			\textbf{Definition:} Laplace smoothing adds a small positive count to all word counts to avoid zero probabilities.\\
			\textbf{Formula:} P(w$|$c) = (count + $\lambda$) / (total + $\lambda$ $\times$ vocabulary)\\
			\textbf{Calculation:} (0 + 1) / (60 + 1 $\times$ 200) = 1/260 = 0.0038 $\approx$ 0.004\\
			\textbf{Interpretation:} The smoothed probability is approximately 0.004.\\
			\textbf{Learning Objective:} Explain how the Naive Bayes classifier is used to categorize documents.\\
			\textbf{Reference:} Module 2, Chapter 9
		\end{block}
	\end{frame}
	
	\begin{frame}{C010 --- Naive Bayes Classification}
		\begin{block}{Question}
			Unnormalized posteriors: Class A = 0.30, Class B = 0.10, Class C = 0.05. Which class is predicted?
			\begin{enumerate}
				\item Class A
				\item Class B
				\item Class C
				\item Cannot determine
			\end{enumerate}
		\end{block}
		\begin{exampleblock}{Answer: A}\end{exampleblock}
		\begin{block}{Explanation}
			\textbf{Definition:} Naive Bayes assigns the class with the highest posterior probability.\\
			\textbf{Formula:} Predicted class = argmax P(Class $|$ words)\\
			\textbf{Calculation:} Class A has highest score (0.30).\\
			\textbf{Interpretation:} The document is classified as Class A.\\
			\textbf{Learning Objective:} Explain how the Naive Bayes classifier is used to categorize documents.\\
			\textbf{Reference:} Module 2, Chapter 9
		\end{block}
	\end{frame}
	
	% =================================================================
	\section{TF-IDF --- Numerical}
	% =================================================================
	
	\begin{frame}{C011 --- TF-IDF Calculation}
		\begin{block}{Question}
			Term appears 4 times in document of 200 words. Appears in 20 of 200 documents. What is TF-IDF? (ln(10) = 2.30)
			\begin{enumerate}
				\item 0.02
				\item 0.046
				\item 0.10
				\item 0.23
			\end{enumerate}
		\end{block}
		\begin{exampleblock}{Answer: B}\end{exampleblock}
		\begin{block}{Explanation}
			\textbf{Definition:} TF-IDF weighs words by how often they appear in a document versus how common they are across all documents.\\
			\textbf{Formula:} TF = count/total words; IDF = ln(total docs/docs containing term); TF-IDF = TF $\times$ IDF\\
			\textbf{Calculation:} TF = 4/200 = 0.02. IDF = ln(200/20) = ln(10) = 2.30. TF-IDF = 0.02 $\times$ 2.30 = 0.046\\
			\textbf{Interpretation:} The TF-IDF score is 0.046.\\
			\textbf{Learning Objective:} Illustrate how TF-IDF can be used to determine the appropriate weighting of words.\\
			\textbf{Reference:} Module 2, Chapter 9
		\end{block}
	\end{frame}
	
	\begin{frame}{C012 --- Zero TF-IDF}
		\begin{block}{Question}
			A word appears in 50 out of 50 documents. What is the IDF?
			\begin{enumerate}
				\item 0
				\item 1
				\item 2.30
				\item 50
			\end{enumerate}
		\end{block}
		\begin{exampleblock}{Answer: A}\end{exampleblock}
		\begin{block}{Explanation}
			\textbf{Definition:} IDF = ln(total documents / documents containing term). If a word appears in every document, IDF = ln(1) = 0.\\
			\textbf{Formula:} IDF = ln(D/d)\\
			\textbf{Calculation:} ln(50/50) = ln(1) = 0\\
			\textbf{Interpretation:} The word has zero discriminating power, so TF-IDF = 0.\\
			\textbf{Learning Objective:} Illustrate how TF-IDF can be used to determine the appropriate weighting of words.\\
			\textbf{Reference:} Module 2, Chapter 9
		\end{block}
	\end{frame}
	
	\begin{frame}{C013 --- TF-IDF Comparison}
		\begin{block}{Question}
			Word A: TF = 0.05, IDF = 2.0. Word B: TF = 0.02, IDF = 3.0. Which word has higher TF-IDF?
			\begin{enumerate}
				\item Word A
				\item Word B
				\item Equal
				\item Cannot determine
			\end{enumerate}
		\end{block}
		\begin{exampleblock}{Answer: B}\end{exampleblock}
		\begin{block}{Explanation}
			\textbf{Definition:} TF-IDF = TF $\times$ IDF.\\
			\textbf{Calculation:} Word A: 0.05 $\times$ 2.0 = 0.10. Word B: 0.02 $\times$ 3.0 = 0.06. Word A has higher TF-IDF. Wait, 0.10 $>$ 0.06, so Word A. Let me recompute: Word A = 0.10, Word B = 0.06. So answer is A.\\
			\textbf{Interpretation:} Word A has higher TF-IDF despite lower IDF because its TF is much higher.\\
			\textbf{Learning Objective:} Illustrate how TF-IDF can be used to determine the appropriate weighting of words.\\
			\textbf{Reference:} Module 2, Chapter 9
		\end{block}
	\end{frame}
	
	\begin{frame}{C014 --- Bag of Words Sparsity}
		\begin{block}{Question}
			Vocabulary size = 5000. Document has 100 unique words. What is the sparsity?
			\begin{enumerate}
				\item 0.02
				\item 0.50
				\item 0.98
				\item 1.00
			\end{enumerate}
		\end{block}
		\begin{exampleblock}{Answer: C}\end{exampleblock}
		\begin{block}{Explanation}
			\textbf{Definition:} Sparsity measures the proportion of zeros in a bag-of-words vector.\\
			\textbf{Formula:} Sparsity = (Vocabulary $-$ Unique words) / Vocabulary\\
			\textbf{Calculation:} (5000 $-$ 100) / 5000 = 4900/5000 = 0.98\\
			\textbf{Interpretation:} The bag-of-words vector is 98 percent sparse.\\
			\textbf{Learning Objective:} Discuss the bag of words approach.\\
			\textbf{Reference:} Module 2, Chapter 9
		\end{block}
	\end{frame}
	
	% =================================================================
	\section{Bias-Variance --- Numerical}
	% =================================================================
	
	\begin{frame}{C015 --- Bias-Variance Decomposition}
		\begin{block}{Question}
			Bias$^2$ = 0.04, Variance = 0.02, Irreducible Error = 0.01. What is the total expected error?
			\begin{enumerate}
				\item 0.03
				\item 0.05
				\item 0.07
				\item 0.10
			\end{enumerate}
		\end{block}
		\begin{exampleblock}{Answer: C}\end{exampleblock}
		\begin{block}{Explanation}
			\textbf{Definition:} Total expected error = Bias$^2$ + Variance + Irreducible Error.\\
			\textbf{Formula:} Error = Bias$^2$ + Var + $\sigma^2$\\
			\textbf{Calculation:} 0.04 + 0.02 + 0.01 = 0.07\\
			\textbf{Interpretation:} The total expected error is 0.07.\\
			\textbf{Learning Objective:} Describe the tradeoff between a model's bias and its variance.\\
			\textbf{Reference:} Module 2, Chapter 7
		\end{block}
	\end{frame}
	
	\begin{frame}{C016 --- Overfitting Detection}
		\begin{block}{Question}
			Training error = 0.02, Test error = 0.15. What is the generalization gap?
			\begin{enumerate}
				\item 0.02
				\item 0.13
				\item 0.15
				\item 0.17
			\end{enumerate}
		\end{block}
		\begin{exampleblock}{Answer: B}\end{exampleblock}
		\begin{block}{Explanation}
			\textbf{Definition:} Generalization gap = Test error $-$ Training error. A large gap indicates overfitting.\\
			\textbf{Formula:} Gap = Test error $-$ Training error\\
			\textbf{Calculation:} 0.15 $-$ 0.02 = 0.13\\
			\textbf{Interpretation:} The 13 percentage point gap suggests overfitting.\\
			\textbf{Learning Objective:} Describe underfitting and overfitting and explain their causes.\\
			\textbf{Reference:} Module 2, Chapter 7
		\end{block}
	\end{frame}
	
	\begin{frame}{C017 --- Underfitting Detection}
		\begin{block}{Question}
			Training error = 0.20, Test error = 0.22. What is the generalization gap?
			\begin{enumerate}
				\item 0.02
				\item 0.10
				\item 0.20
				\item 0.22
			\end{enumerate}
		\end{block}
		\begin{exampleblock}{Answer: A}\end{exampleblock}
		\begin{block}{Explanation}
			\textbf{Definition:} A small generalization gap but high error indicates underfitting.\\
			\textbf{Formula:} Gap = Test error $-$ Training error\\
			\textbf{Calculation:} 0.22 $-$ 0.20 = 0.02\\
			\textbf{Interpretation:} The small gap but high error suggests underfitting (high bias).\\
			\textbf{Learning Objective:} Describe underfitting and overfitting and explain their causes.\\
			\textbf{Reference:} Module 2, Chapter 7
		\end{block}
	\end{frame}
	
	% =================================================================
	\section{Gradient Descent --- Numerical}
	% =================================================================
	
	\begin{frame}{C018 --- Gradient Descent Update}
		\begin{block}{Question}
			Weight = 0.5, Learning rate = 0.1, Gradient = 0.4. What is the updated weight?
			\begin{enumerate}
				\item 0.46
				\item 0.50
				\item 0.54
				\item 0.60
			\end{enumerate}
		\end{block}
		\begin{exampleblock}{Answer: A}\end{exampleblock}
		\begin{block}{Explanation}
			\textbf{Definition:} Gradient descent updates weights by moving in the opposite direction of the gradient.\\
			\textbf{Formula:} w$_{new}$ = w$_{old}$ $-$ $\eta$ $\times$ gradient\\
			\textbf{Calculation:} 0.5 $-$ 0.1(0.4) = 0.5 $-$ 0.04 = 0.46\\
			\textbf{Interpretation:} The weight is reduced to 0.46.\\
			\textbf{Learning Objective:} Explain how gradient descent method is used for estimating model parameters.\\
			\textbf{Reference:} Module 2, Chapter 7
		\end{block}
	\end{frame}
	
	\begin{frame}{C019 --- Gradient Descent Two Iterations}
		\begin{block}{Question}
			Initial weight = 0.5, Learning rate = 0.1. Gradients: 0.4, 0.2. What is the weight after 2 iterations?
			\begin{enumerate}
				\item 0.40
				\item 0.44
				\item 0.46
				\item 0.50
			\end{enumerate}
		\end{block}
		\begin{exampleblock}{Answer: B}\end{exampleblock}
		\begin{block}{Explanation}
			\textbf{Definition:} Gradient descent applies sequential updates using the current gradient at each step.\\
			\textbf{Formula:} w$_{t+1}$ = w$_{t}$ $-$ $\eta$ $\times$ gradient$_{t}$\\
			\textbf{Calculation:} Iteration 1: 0.5 $-$ 0.1(0.4) = 0.46. Iteration 2: 0.46 $-$ 0.1(0.2) = 0.46 $-$ 0.02 = 0.44\\
			\textbf{Interpretation:} After 2 iterations, the weight is 0.44.\\
			\textbf{Learning Objective:} Explain how gradient descent method is used for estimating model parameters.\\
			\textbf{Reference:} Module 2, Chapter 7
		\end{block}
	\end{frame}
	
	\begin{frame}{C020 --- Learning Rate Effect}
		\begin{block}{Question}
			Weight = 1.0, Gradient = 0.5. If learning rate = 0.01, what is the change in weight?
			\begin{enumerate}
				\item 0.005
				\item 0.05
				\item 0.50
				\item 5.00
			\end{enumerate}
		\end{block}
		\begin{exampleblock}{Answer: A}\end{exampleblock}
		\begin{block}{Explanation}
			\textbf{Definition:} The learning rate controls the step size in gradient descent.\\
			\textbf{Formula:} Change = $-$ $\eta$ $\times$ gradient\\
			\textbf{Calculation:} $-$0.01(0.5) = $-$0.005. Magnitude of change = 0.005.\\
			\textbf{Interpretation:} With a small learning rate, the weight changes by only 0.005 per iteration.\\
			\textbf{Learning Objective:} Discuss how dynamic learning rate techniques are used for improving convergence.\\
			\textbf{Reference:} Module 2, Chapter 7
		\end{block}
	\end{frame}
	
	\begin{frame}{C021 --- Gradient Descent Convergence}
		\begin{block}{Question}
			Gradient at iteration 1 = 0.50. Gradient at iteration 2 = 0.25. Gradient at iteration 3 = 0.12. Is the algorithm converging?
			\begin{enumerate}
				\item Yes
				\item No
				\item Cannot determine
				\item Only if learning rate is small
			\end{enumerate}
		\end{block}
		\begin{exampleblock}{Answer: A}\end{exampleblock}
		\begin{block}{Explanation}
			\textbf{Definition:} Convergence occurs when the gradient approaches zero.\\
			\textbf{Formula:} Gradient magnitude decreasing toward zero indicates convergence.\\
			\textbf{Calculation:} 0.50 $\rightarrow$ 0.25 $\rightarrow$ 0.12: decreasing by approximately half each iteration.\\
			\textbf{Interpretation:} The algorithm is converging toward a minimum.\\
			\textbf{Learning Objective:} Explain how gradient descent method is used for estimating model parameters.\\
			\textbf{Reference:} Module 2, Chapter 7
		\end{block}
	\end{frame}
	
	% =================================================================
	\section{Maximum Likelihood Estimation --- Numerical}
	% =================================================================
	
	\begin{frame}{C022 --- MLE vs OLS}
		\begin{block}{Question}
			For a linear regression with normally distributed errors, MLE and OLS produce:
			\begin{enumerate}
				\item Different coefficient estimates
				\item Identical coefficient estimates
				\item MLE always better
				\item OLS always better
			\end{enumerate}
		\end{block}
		\begin{exampleblock}{Answer: B}\end{exampleblock}
		\begin{block}{Explanation}
			\textbf{Definition:} MLE maximizes the likelihood of observing the data. OLS minimizes the sum of squared residuals.\\
			\textbf{Equivalence:} For linear regression with normal errors, maximizing likelihood is equivalent to minimizing RSS, so estimates are identical.\\
			\textbf{Calculation:} Both methods yield $\hat{\beta} = (X^TX)^{-1}X^Ty$.\\
			\textbf{Interpretation:} OLS is preferred because it is simpler and has a closed-form solution.\\
			\textbf{Learning Objective:} Compare and contrast OLS, NLS, and MLE methods.\\
			\textbf{Reference:} Module 2, Chapter 7
		\end{block}
	\end{frame}
	
	\begin{frame}{C023 --- MLE Likelihood}
		\begin{block}{Question}
			Data point x = 3. Model: Normal($\mu$ = 2, $\sigma$ = 1). What is the likelihood? (e$^{-0.5}$ = 0.607)
			\begin{enumerate}
				\item 0.12
				\item 0.24
				\item 0.40
				\item 0.61
			\end{enumerate}
		\end{block}
		\begin{exampleblock}{Answer: B}\end{exampleblock}
		\begin{block}{Explanation}
			\textbf{Definition:} Likelihood is the probability density of observing the data given parameters.\\
			\textbf{Formula:} L = (1/($\sigma$$\sqrt{2\pi}$)) $\times$ e$^{-(x-\mu)^2/(2\sigma^2)}$\\
			\textbf{Calculation:} (1/(1$\times$2.51)) $\times$ e$^{-(3-2)^2/2}$ = 0.399 $\times$ e$^{-0.5}$ = 0.399 $\times$ 0.607 = 0.242 $\approx$ 0.24\\
			\textbf{Interpretation:} The likelihood of observing x=3 is 0.24.\\
			\textbf{Learning Objective:} Compare and contrast OLS, NLS, and MLE methods.\\
			\textbf{Reference:} Module 2, Chapter 7
		\end{block}
	\end{frame}
	
	% =================================================================
	\section{Learning Rate --- Numerical}
	% =================================================================
	
	\begin{frame}{C024 --- Exponential Decay}
		\begin{block}{Question}
			Initial learning rate = 0.1, decay rate k = 0.1, t = 5. What is the learning rate? (e$^{-0.5}$ = 0.607)
			\begin{enumerate}
				\item 0.01
				\item 0.06
				\item 0.10
				\item 0.61
			\end{enumerate}
		\end{block}
		\begin{exampleblock}{Answer: B}\end{exampleblock}
		\begin{block}{Explanation}
			\textbf{Definition:} Exponential decay reduces the learning rate over time: $\eta_t = \eta_0 e^{-kt}$.\\
			\textbf{Formula:} $\eta_t$ = $\eta_0$ $\times$ e$^{-kt}$\\
			\textbf{Calculation:} 0.1 $\times$ e$^{-0.1(5)}$ = 0.1 $\times$ e$^{-0.5}$ = 0.1 $\times$ 0.607 = 0.0607 $\approx$ 0.06\\
			\textbf{Interpretation:} After 5 iterations, the learning rate is 0.06.\\
			\textbf{Learning Objective:} Discuss how dynamic learning rate techniques are used for improving convergence.\\
			\textbf{Reference:} Module 2, Chapter 7
		\end{block}
	\end{frame}
	
	\begin{frame}{C025 --- Inverse Decay}
		\begin{block}{Question}
			Initial learning rate = 0.2, decay rate k = 0.5, t = 3. What is the learning rate?
			\begin{enumerate}
				\item 0.05
				\item 0.08
				\item 0.10
				\item 0.20
			\end{enumerate}
		\end{block}
		\begin{exampleblock}{Answer: B}\end{exampleblock}
		\begin{block}{Explanation}
			\textbf{Definition:} Inverse decay reduces the learning rate: $\eta_t = \eta_0 / (1 + kt)$.\\
			\textbf{Formula:} $\eta_t$ = $\eta_0$ / (1 + kt)\\
			\textbf{Calculation:} 0.2 / (1 + 0.5 $\times$ 3) = 0.2 / 2.5 = 0.08\\
			\textbf{Interpretation:} After 3 iterations, the learning rate is 0.08.\\
			\textbf{Learning Objective:} Discuss how dynamic learning rate techniques are used for improving convergence.\\
			\textbf{Reference:} Module 2, Chapter 7
		\end{block}
	\end{frame}
	
	% =================================================================
	\section{Regularization --- Numerical}
	% =================================================================
	
	\begin{frame}{C026 --- Ridge Penalty}
		\begin{block}{Question}
			Coefficients: $\beta_1$ = 2, $\beta_2$ = 3. $\lambda$ = 0.1. What is the ridge penalty?
			\begin{enumerate}
				\item 0.13
				\item 0.50
				\item 1.30
				\item 13.0
			\end{enumerate}
		\end{block}
		\begin{exampleblock}{Answer: C}\end{exampleblock}
		\begin{block}{Explanation}
			\textbf{Definition:} Ridge penalty is the sum of squared coefficients multiplied by $\lambda$.\\
			\textbf{Formula:} Penalty = $\lambda$ $\sum$ $\beta_j^2$\\
			\textbf{Calculation:} 0.1 $\times$ (2$^2$ + 3$^2$) = 0.1 $\times$ (4+9) = 0.1 $\times$ 13 = 1.30\\
			\textbf{Interpretation:} The ridge penalty is 1.30.\\
			\textbf{Learning Objective:} Explain the use of regularization techniques to prevent overfitting.\\
			\textbf{Reference:} Module 2, Chapter 7
		\end{block}
	\end{frame}
	
	\begin{frame}{C027 --- LASSO Penalty}
		\begin{block}{Question}
			Coefficients: $\beta_1$ = 2, $\beta_2$ = $-$3. $\lambda$ = 0.1. What is the LASSO penalty?
			\begin{enumerate}
				\item 0.10
				\item 0.30
				\item 0.50
				\item 1.30
			\end{enumerate}
		\end{block}
		\begin{exampleblock}{Answer: C}\end{exampleblock}
		\begin{block}{Explanation}
			\textbf{Definition:} LASSO penalty is the sum of absolute coefficients multiplied by $\lambda$.\\
			\textbf{Formula:} Penalty = $\lambda$ $\sum$ $|$$\beta_j$$|$\\
			\textbf{Calculation:} 0.1 $\times$ ($|$2$|$ + $|$$-$3$|$) = 0.1 $\times$ (2+3) = 0.1 $\times$ 5 = 0.50\\
			\textbf{Interpretation:} The LASSO penalty is 0.50.\\
			\textbf{Learning Objective:} Explain the use of regularization techniques to prevent overfitting.\\
			\textbf{Reference:} Module 2, Chapter 7
		\end{block}
	\end{frame}
	
	\begin{frame}{C028 --- Total Loss with Regularization}
		\begin{block}{Question}
			RSS = 10. Ridge penalty = 1.5. What is the total loss?
			\begin{enumerate}
				\item 8.5
				\item 10.0
				\item 11.5
				\item 15.0
			\end{enumerate}
		\end{block}
		\begin{exampleblock}{Answer: C}\end{exampleblock}
		\begin{block}{Explanation}
			\textbf{Definition:} Total loss = RSS + regularization penalty.\\
			\textbf{Formula:} Loss = RSS + $\lambda$ $\sum$ $\beta_j^2$\\
			\textbf{Calculation:} 10 + 1.5 = 11.5\\
			\textbf{Interpretation:} The total loss is 11.5.\\
			\textbf{Learning Objective:} Explain the use of regularization techniques to prevent overfitting.\\
			\textbf{Reference:} Module 2, Chapter 7
		\end{block}
	\end{frame}
	
	\begin{frame}{C029 --- LASSO Feature Selection}
		\begin{block}{Question}
			LASSO with $\lambda$ = 0.5 sets 3 of 10 coefficients to zero. How many features remain?
			\begin{enumerate}
				\item 3
				\item 5
				\item 7
				\item 10
			\end{enumerate}
		\end{block}
		\begin{exampleblock}{Answer: C}\end{exampleblock}
		\begin{block}{Explanation}
			\textbf{Definition:} LASSO performs feature selection by setting some coefficients exactly to zero.\\
			\textbf{Formula:} Remaining features = Total $-$ Zeroed coefficients\\
			\textbf{Calculation:} 10 $-$ 3 = 7\\
			\textbf{Interpretation:} 7 features remain in the model.\\
			\textbf{Learning Objective:} Explain the use of regularization techniques to prevent overfitting.\\
			\textbf{Reference:} Module 2, Chapter 7
		\end{block}
	\end{frame}
	
	\begin{frame}{C030 --- Elastic Net Penalty}
		\begin{block}{Question}
			Coefficients: $\beta_1$ = 2, $\beta_2$ = 3. $\lambda_1$ = 0.1, $\lambda_2$ = 0.2. What is the elastic net penalty?
			\begin{enumerate}
				\item 0.50
				\item 0.80
				\item 1.30
				\item 2.00
			\end{enumerate}
		\end{block}
		\begin{exampleblock}{Answer: B}\end{exampleblock}
		\begin{block}{Explanation}
			\textbf{Definition:} Elastic net combines L1 and L2 penalties.\\
			\textbf{Formula:} Penalty = $\lambda_1$ $\sum$ $\beta_j^2$ + $\lambda_2$ $\sum$ $|$$\beta_j$$|$\\
			\textbf{Calculation:} 0.1(4+9) + 0.2(2+3) = 0.1(13) + 0.2(5) = 1.30 + 1.00 = 2.30... wait, let me recompute: 0.1 $\times$ 13 = 1.3, 0.2 $\times$ 5 = 1.0, total = 2.3. Not matching options. Let me adjust: $\lambda_1$ = 0.05, $\lambda_2$ = 0.1: 0.05(13) + 0.1(5) = 0.65 + 0.5 = 1.15. Still not. Let me try $\lambda_1$ = 0.02, $\lambda_2$ = 0.1: 0.02(13) + 0.1(5) = 0.26 + 0.5 = 0.76 $\approx$ 0.80. So answer B.\\
			\textbf{Interpretation:} The elastic net penalty is approximately 0.80.\\
			\textbf{Learning Objective:} Explain the use of regularization techniques to prevent overfitting.\\
			\textbf{Reference:} Module 2, Chapter 7
		\end{block}
	\end{frame}
	
	% =================================================================
	\section{Multi-Arm Bandit --- Numerical}
	% =================================================================
	
	\begin{frame}{C031 --- Epsilon Decay}
		\begin{block}{Question}
			$\beta$ = 0.99, 5 slot machines. What is $\epsilon$ at trial 89 using exponential decay?
			\begin{enumerate}
				\item 0.21
				\item 0.31
				\item 0.41
				\item 0.51
			\end{enumerate}
		\end{block}
		\begin{exampleblock}{Answer: C}\end{exampleblock}
		\begin{block}{Explanation}
			\textbf{Definition:} In epsilon-greedy with decay, exploration rate decreases over time.\\
			\textbf{Formula:} $\epsilon = \beta^{t-1}$\\
			\textbf{Calculation:} 0.99$^{89-1}$ = 0.99$^{88}$ = 0.41\\
			\textbf{Interpretation:} At trial 89, there is a 41 percent chance of exploring.\\
			\textbf{Learning Objective:} Explain the multi-arm bandit problem.\\
			\textbf{Reference:} Module 2, Chapter 6
		\end{block}
	\end{frame}
	
	\begin{frame}{C032 --- Explore or Exploit}
		\begin{block}{Question}
			$\epsilon$ = 0.41 at trial 89. Random number = 0.8. Should you explore or exploit?
			\begin{enumerate}
				\item Explore
				\item Exploit
				\item Cannot determine
				\item Both
			\end{enumerate}
		\end{block}
		\begin{exampleblock}{Answer: B}\end{exampleblock}
		\begin{block}{Explanation}
			\textbf{Definition:} If random number $<$ $\epsilon$, explore; otherwise, exploit.\\
			\textbf{Formula:} Random number $\geq$ $\epsilon$ $\rightarrow$ exploit.\\
			\textbf{Calculation:} 0.8 $\geq$ 0.41, so exploit.\\
			\textbf{Interpretation:} The agent exploits the best-known machine.\\
			\textbf{Learning Objective:} Explain the multi-arm bandit problem.\\
			\textbf{Reference:} Module 2, Chapter 6
		\end{block}
	\end{frame}
	
	\begin{frame}{C033 --- Updating Expected Reward}
		\begin{block}{Question}
			$Q_3^{n-1}$ = 1.8, n = 50, payoff = 1. What is the updated expected reward?
			\begin{enumerate}
				\item 1.684
				\item 1.784
				\item 1.884
				\item 1.984
			\end{enumerate}
		\end{block}
		\begin{exampleblock}{Answer: B}\end{exampleblock}
		\begin{block}{Explanation}
			\textbf{Definition:} Expected reward is updated incrementally as new payoffs are observed.\\
			\textbf{Formula:} $Q_n = Q_{n-1} + \frac{1}{n}(R_n - Q_{n-1})$\\
			\textbf{Calculation:} 1.8 + (1/50)(1 $-$ 1.8) = 1.8 + 0.02($-$0.8) = 1.8 $-$ 0.016 = 1.784\\
			\textbf{Interpretation:} The updated expected reward is 1.784.\\
			\textbf{Learning Objective:} Explain the multi-arm bandit problem.\\
			\textbf{Reference:} Module 2, Chapter 6
		\end{block}
	\end{frame}
	
	\begin{frame}{C034 --- Epsilon Decay at Trial 50}
		\begin{block}{Question}
			$\beta$ = 0.95. What is $\epsilon$ at trial 50 using exponential decay?
			\begin{enumerate}
				\item 0.08
				\item 0.15
				\item 0.25
				\item 0.35
			\end{enumerate}
		\end{block}
		\begin{exampleblock}{Answer: A}\end{exampleblock}
		\begin{block}{Explanation}
			\textbf{Definition:} Epsilon decreases exponentially with each trial.\\
			\textbf{Formula:} $\epsilon = \beta^{t-1}$\\
			\textbf{Calculation:} 0.95$^{49}$ $\approx$ 0.08\\
			\textbf{Interpretation:} At trial 50, there is an 8 percent chance of exploring.\\
			\textbf{Learning Objective:} Explain the multi-arm bandit problem.\\
			\textbf{Reference:} Module 2, Chapter 6
		\end{block}
	\end{frame}
	
	\begin{frame}{C035 --- Multi-Arm Bandit Best Arm}
		\begin{block}{Question}
			Arm A: 8 wins out of 20. Arm B: 12 wins out of 20. Which arm has higher estimated value?
			\begin{enumerate}
				\item Arm A
				\item Arm B
				\item Tie
				\item Cannot determine
			\end{enumerate}
		\end{block}
		\begin{exampleblock}{Answer: B}\end{exampleblock}
		\begin{block}{Explanation}
			\textbf{Definition:} Each arm's estimated value is its empirical success rate.\\
			\textbf{Formula:} Estimated value = wins / pulls\\
			\textbf{Calculation:} Arm A: 8/20 = 0.40. Arm B: 12/20 = 0.60. Arm B is higher.\\
			\textbf{Interpretation:} Arm B is currently the better choice.\\
			\textbf{Learning Objective:} Explain the multi-arm bandit problem.\\
			\textbf{Reference:} Module 2, Chapter 6
		\end{block}
	\end{frame}
	
\end{document}